\section{Seguridad, Pruebas de Rendimiento y Despliegue con Docker}
\subsection{Auditoría de seguridad OWASP}

Se realizó la evaluación de seguridad sobre la aplicación usando el OWASP Top 10:2021 (A01--A07) más XSS. La
Tabla~\ref{tab:owasp} resume la vulnerabilidad, el vector de ataque, la contramedida implementada y la evidencia.

\begin{table}[H]
\centering
\footnotesize
\setlength{\tabcolsep}{4pt}
\begin{tabular}{@{}p{2.4cm}p{2.9cm}p{3.3cm}p{2.6cm}@{}}
\toprule
\textbf{Vulnerabilidad} & \textbf{Vector de ataque} & \textbf{Contramedida implementada} & \textbf{Evidencia} \\
\midrule
A01 Control de acceso & Acceso no autorizado a /api/solicitudes & Reglas @PreAuthorize y SecurityConfig (ADMIN, DOCENTE, COORDINADOR) & OWASP-AUDIT.md \\
A02 Fallo criptográfico & Robo de credenciales/token & BCryptPasswordEncoder, JWT HS256, cookie HttpOnly & SecurityConfig + AuthController \\
A03 Inyección & SQL/JPQL malicioso en solicitudes & Validación Jakarta @Valid y Spring Data JPA parametrizado & OWASP-AUDIT.md \\
A05 Malconfiguración & Cabeceras inseguras & Desactivación CSRF adecuada para Stateless JWT, CORS acotado & SecurityConfig.java \\
A07 Fallos de autenticación & Fuerza bruta en login & Autenticación gestionada por AuthenticationManager & AuthController.java \\
XSS & Script malicioso en el navegador & Saneamiento automático de plantillas Angular 21 & Frontend Angular \\
\bottomrule
\end{tabular}
\caption{Resumen de la auditoría OWASP aplicada al PFC Pre-Sustentaciones.}
\label{tab:owasp}
\end{table}

\subsection{Prueba de carga con k6}

La metodología de prueba de carga del PFC se realizó con \textbf{k6} contra endpoints representativos. Se ejecutaron
3 corridas independientes (`run1`, `run2`, `run3`). Los resultados confirmaron baja latencia y alta capacidad de
procesamiento bajo carga simulada, registrando métricas crudas en `docs/pruebas/k6/`.

\subsection{Containerización con Docker Compose}

Docker Compose permite desplegar la infraestructura de la aplicación con aislamiento de dependencias. El
\texttt{docker-compose.yml} del PFC define los servicios:
\begin{itemize}
    \item \textbf{postgres} (postgres:15-alpine, puerto 5432:5432, volumen \texttt{postgres\_data}).
    \item \textbf{redis} (redis:7-alpine, puerto 6379:6379) con healthcheck \texttt{redis-cli ping}.
\end{itemize}

El despliegue se ejecuta de forma unificada mediante `docker compose up -d` o con el `Makefile` (`make up`).
